\documentclass[11pt]{article}
\usepackage[margin=1in]{geometry}

\usepackage{graphicx} 
\usepackage{kotex}
\usepackage{amsmath, amssymb, amsthm}
\usepackage{mathtools}
\usepackage{enumitem}
\usepackage{xcolor}
\usepackage{hyperref}
\usepackage[capitalize,nameinlink]{cleveref}

\hypersetup{
    colorlinks=true,
    linkcolor=blue,
    citecolor=blue,
    urlcolor=blue
}

\let\oldemptyset\emptyset
\let\emptyset\varnothing

\newtheorem{theorem}{Theorem}[section]
\newtheorem{lemma}[theorem]{Lemma}
\newtheorem{proposition}[theorem]{Proposition}
\newtheorem{corollary}[theorem]{Corollary}

\theoremstyle{definition}
\newtheorem{definition}[theorem]{Definition}
\newtheorem{example}[theorem]{Example}

\theoremstyle{remark}
\newtheorem{remark}[theorem]{Remark}

\DeclareMathOperator*{\argmax}{arg\,max}
\DeclareMathOperator*{\argmin}{arg\,min}

\newcommand{\R}{\mathbb{R}}
\newcommand{\C}{\mathbb{C}}
\newcommand{\Z}{\mathbb{Z}}
\newcommand{\Q}{\mathbb{Q}}
\newcommand{\N}{\mathbb{N}}

\newcommand{\eps}{\varepsilon}
\newcommand{\dd}{\,\mathrm{d}}
\newcommand{\id}{\mathrm{id}}
\newcommand{\im}{\operatorname{im}}
\newcommand{\rank}{\operatorname{rank}}
\newcommand{\Hom}{\operatorname{Hom}}
\newcommand{\ntoinf}{n \to \infty}

\newcommand{\inner}[2]{\left\langle #1, #2 \right\rangle}
\DeclarePairedDelimiter{\abs}{\lvert}{\rvert}
\DeclarePairedDelimiter{\norm}{\lVert}{\rVert}
\DeclarePairedDelimiter{\innerp}{\langle}{\rangle}

\title{해개연1 과제2}
\author{윤승현}
\date{2026.04}

\begin{document}

\maketitle

\subsection*{5.}

\begin{lemma}
    Let $b_n=\sup_{m \geq n} a_m$, then
    \[
    \limsup_{\ntoinf} a_n = \lim_{\ntoinf} b_n
    \]
    \label{lemma:limsup}
\end{lemma}

\begin{proof}
    The $b_n$ is monotically decreasing sequence. ($\because$ $\{a_k\}_{k \geq n+1} \subset \{a_k\}_{k \geq n}$) In extended real, $L:=\lim b_n=\inf b_n$. 

    Set a subsequence $\{a_{n_j}\}$ converges to $c$ s.t. $n_j > N$, $a_{n_j} \leq b_N$ since $a_{n_j} \in \{a_k\}_{k \geq N}$. $c \leq b_N$, then $c \leq \inf b_n = L$ Then $\limsup a_n \leq L$.\\
    For sufficiently large $k$, $b_n < L + 1/k$ for $n \geq N_k$. Since $b_n = \sup_{m \geq N_k} a_m$, $\exists n_k > N_k > k$,
    
    \[
        a_{n_k}>b_{N_k}-1/k \implies a_{n_k}>L-1/k
    \]

    \[
        0 \leq d(a_{n_k}, L) < 1/k
    \]

    Then $L \leq \limsup a_n$. Therefore $\limsup a_n=\lim b_n$
\end{proof}

For $n$, $\sup_{m \geq n}(a_m+b_m) \leq \sup(a_m) + \sup(b_m)$. ($\because \, a_m \leq \sup_{m \geq n}(a_m)$ similarly $b_m$ then its sum is bounded) With \cref{lemma:limsup}, we have $\limsup(a_n+b_n)\leq\lim(\sup(a_m) + \sup(b_m))=\limsup(a_n)+\limsup(b_n)$ ($\because$ it is not form of $\infty - \infty$)

\subsection*{8.}

Since $b_n$ is monotone and bounded. For the monotically decreasing case, it converges to $\alpha$ by Thm 3.14. Let $c_n=b_n-\alpha$. Then $c_0 \geq c_1 \geq ...$, $\lim b_n=0$ by Thm 3.3. Similarly, for the monotically increasing case, it converges to $\beta$ and let $c_n=-b_n+\beta$ satisfying previous statements.\\
By Thm 3.22 (Cauchy criterion), $\eps>0$ $\exists N$ s.t. 

\[
\abs*{\sum_{n}^{m} a_n} < \eps
\]

partial sums $A_n$ of $\sum a_n$ form a bounded sequence. (for $n \geq N$)\\
By Thm 3.42, $\sum a_n b_n$ converges.

\subsection*{14.}

\subsubsection*{a}

For $\eps > 0$, $\exists N$ s.t. $d(s_n,s)<\eps$ for $n \geq N$. ($\because$ $s_n$ converges)

\[
\sigma_n = \frac{\sigma_N+(s_{N+1}+s_{N+2}+...+s_n)}{n+1}
\]

By Thm 3.20, $|\frac{\sigma_N-sN}{n+1}| \to 0$. Then $\exists N_1$ s.t.

\[
|\frac{\sigma_N-sN}{n+1}|<\frac{N+1}{n+1}\eps
\]

for $n \geq N_1$. Then, 

\[
|\sigma_n-s| < |\frac{\sigma_N-sN}{n+1}|+\frac{(n-N)}{n+1}\eps < \eps
\]

for $n \geq max(N, N_1)$.

\subsubsection*{b}

Let $s_n=(-1)^n$. It does not converge, since $\limsup s_n=1$ and $\liminf s_n=-1$.

But $\inf_n \sigma_m = 0$ and

\[
\sup_n \sigma_m=\begin{cases}
    1/(n+1) & \text{if n is even}\\
    1/(n+2) & \text{otherwise}
\end{cases}
\]

Since $\limsup \sigma_n=0$ and $\liminf \sigma_n=0$, by \cref{lemma:limsup}, $\lim \sigma_n=0$

\subsubsection*{c}

Let

\[
s_n = \begin{cases}
    2^k & n=k\times 2^k \\
    2^{-n} & \text{otherwise}
\end{cases}
\]

It's clearly $s_n>0$ and $\sup_{k^2 2^k} s_m = 2^k$, then by \cref{lemma:limsup}, $\limsup s_m = \infty$.

\[
    \frac{\sum_{k=0}^{n} 2^{-k}}{n} < \frac{2(1-2^{-n-1})}{n}
\]

and

\[
\frac{2^1+2^2+...+2^k}{k\times2^k} < 2/k
\]

Then,

\[
0 < \sigma_n < \frac{\sum_{k=0}^{n} 2^{-k}}{n} + \frac{2^1+2^2+...+2^k}{k\times2^k} < 4/n
\]

By Thm 3.25, $\sigma_n$ converges to $0$.

\subsubsection*{d}

\[
\frac{1}{n+1}\sum_{1}^{n} ka_k =\frac{(n+1)s_n-(s_n+s_{n-1}+...)}{n+1}=s_n-\sigma_n
\]

Set $g_n=na_n$. From (a), $\lim g_n=0$ then $\lim \sigma_{g_n}=0$. With 

\[
\sigma_{g_n}=\frac{1}{n+1}\sum_{1}^{n} ka_k=s_n-\sigma_n
\]

and by thm 3.3, $\lim s_n = \lim (s_n-\sigma_n) + \lim \sigma_n = \lim \sigma_n$.

\subsubsection*{e}


\subsection*{17.}


(a) and (b)

\[
x_{n+2}=\frac{2\alpha+(\alpha+1)x_n}{1+\alpha+2x_n}
\]

For $x_n < \sqrt{\alpha}$, $x_{n+2} > x_n$, otherwise $x_n > \sqrt{\alpha}$ then $x_{n+2} < x_n$. And from following, $x_n > \sqrt{\alpha}$

\[
\frac{\alpha-x_n^2}{1+x_n} < \frac{\alpha-x_n^2}{\sqrt{\alpha}+x_n} = \sqrt{\alpha}-x_n
\]

$x_{n+1} < \sqrt{\alpha}$. Similarly, $x_n < \sqrt{\alpha}$ then $x_{n+1} > \sqrt{\alpha}$.\\
Then $x_{2n-1}>\sqrt{\alpha}$ and $x_{2n}<\sqrt{\alpha}$ with monotonically decreasing and increasing.\\

(c) Let

\[
(x^*)_n = \sup_{m \geq n} x_m,\; (x_*)_n = \inf_{m \geq n} x_m
\]

$(x*)_n=x_{2*\lfloor{(n+1)/2}\rfloor-1}$. This is monotically decreasing sequence with bounded below. ($\because \, x_{2n-1}>\sqrt{\alpha}$) By Thm 3.14, it converges.\\

(d)

\subsection*{21.}

Choose $p_n \in E_n$. Then $\{p_n\}$ is Cauchy sequence. ($\because$ $p_n, p_{n+1}... \in E'_n$  $E'_n \subset E_n \implies \text{diam} E'_n \leq \text{diam} E_n$) From Def 3.12, every Cauchy sequence in complete metric space converges, $p_n \to p \in X$. Suppose some $k$, $p$ is not in $E_k$. Since it converges, $\eps > 0$ then $\exists N$ s.t. $d(p,p_n) < \eps$ for $n \geq N$, which implies $N_\eps(p) \cap E_k \neq \emptyset$. Since $E_k$ is closed, $p \in E_k$.\\
Suppose $q \in \bigcap E_n$ s.t. $p \neq q$. Then $\text{diam} E_n \geq d(p,q)$ contradicts to hypothesis. Therefore it consists of exactly one point.

\subsection*{22.}

Let an open ball $C_1$ s.t. $\text{diam} \overline{C_1} < 1$ and $\overline{C_1} \subset G_1$. Assume $C_n \cap G_n$ is non-empty. There's an open ball $C_{n+1}$ with $\text{diam} \overline{C_{n+1}} < 1/n$ and $\overline{C_{n+1}} \subset C_n \cap G_n$. ($\because \, C_n, G_n$ is open and finite intersection also be open) Since $G_{n+1}$ is dense, $C_{n+1} \cap G_{n+1}$ is non-empty.

Then set $E_n=\overline{C_n}$ is closed and $E_1 \supset E_2 \supset E_3 \supset ...$. From Exercies 21. , there's at least one point in $\bigcap G_n$.

\end{document}